% ID: SEL-D0049-OA-E04
% ESTADO: cotejada. Segunda pasada visual y borrador compilado revisados.
% Fuente: Selectividad/sel_2013_matematicas_aplicadas/junio/reserva b/EXAMEN 1-ANDALUCÍA.pdf
% SHA-256: ec6fd00d0daee547aa6128f3f6e1099278ee551f7f80917b331a1772764317b3
% Página: 1; opción A; ejercicio original 4.

Queremos estudiar la proporción de personas de una población que acceden a internet a través de teléfono móvil. Para ello hacemos una encuesta a una muestra aleatoria de 400 personas de esa población, y obtenemos que 240 de ellas acceden a internet a través del móvil.
\begin{enumerate}[label=\alph*)]
\item \textbf{(1.75 puntos)} Determine un intervalo de confianza, al 98.5\%, para la proporción de personas de esa población que acceden a internet a través del teléfono móvil.
\item \textbf{(0.75 puntos)} Razone el efecto que tendría sobre la amplitud del intervalo de confianza el aumento o disminución del tamaño de la muestra, suponiendo que se mantuvieran la misma proporción muestral y el mismo nivel de confianza.
\end{enumerate}
